\begin{theorem}[\textbf{формула коплощади}]
	Пусть $f : \R^{n \geq 2} \to \R$ интегрируема по Лебегу, тогда для почти всех $r \in \R_+$ функция $f$ интегрируема по сфере $S_r(0) = \{x \in \R^n : |x| = r\}$, причем справедлива формула:
	\[ \int_{\R^n} f(x) \, d\mu(x) = \int_0^{+\infty} \left( \int_{S_r(0)} f(x) \, d\nu (x) \right) \, d \mu(r). \]
\end{theorem}

\begin{myproof}
	Пусть $\R^n = H_+ \cup H_- \cup \{x_n = 0\}$, где $H_{\pm} = \{x \in \R^n : \pm x_n > 0\}$, и обозначим $x = (y, x_n)$, $y = (x_1, \dotsc, x_{n-1})$.
	Вычислим интеграл по $H_+$. 
	Рассмотрим отображение $\phi : B_1(0) \times \R_+ \to H_+$ по правилу $\phi(y, r) = (ry, \, r \sqrt{1 - |y|^2})$.
	Это отображение взаимно однозначно, гладкое и имеет гладкое обратное: $\phi^{-1}(x) = (y/ |x|, |x|)$ (в этом можно убедиться, разрешив систему $u = ry$, $v = r \sqrt{1 - |y|^2}$, где $x = (u, v)$). 
	Следовательно, $\phi$ является диффеоморфизмом. 
	Найдем его матрицу Якоби, она имеет следующий вид: 
	\[ D \phi(y, r) = \left( \begin{array}{cccc|c}
		r & 0 & \dots & 0 & x_1 \\
		0 & r & \dots & 0 & x_2 \\
		\vdots & \vdots & \ddots & \vdots & \vdots \\
		0 & 0 & \dots & r & x_{n-1} \\
		\hline
		\frac{-r x_1}{\sqrt{1 - |y|^2}} & \frac{-r x_2}{\sqrt{1 - |y|^2}} & \dots & \frac{-r x_{n-1}}{\sqrt{1 - |y|^2}} & \sqrt{1 - |y|^2}
	\end{array} \right). \]
	Найдем определитель: вынесем из первых $n-1$ столбцов $r$ и прибавим к последней строке комбинацию верхних строк так, чтобы обнулить левый нижний блок. 
	Тогда $\det D \phi = r^{n-1} \left( \sqrt{1 - |y|^2} + \frac{x_1^2 + \dotsc + x^2_{n-1}}{\sqrt{1 - |y|^2}} \right) = \frac{r^{n-1}}{\sqrt{1 - |y|^2}}$.
	По формуле замены и теореме Фубини имеем: 
	\[ \int_{H_+} f \, d\mu = \int_{0}^{+\infty} \left( \int_{B_1(0)} f (ry, r \sqrt{1 - |y|^2}) \cdot \tfrac{r^{n-1}}{\sqrt{1 - |y|^2}} \, dy  \right) \, dr. \]
	Внутренний интеграл равен поверхностному интегралу по полусфере $S_r(0) \cap H_+$ по предыдущему примеру.
	Аналогичное равенство получим для $H_-$. Складывая полученные равенства, и учитывая, что $\{x_n = 0\}$ имеет меру нуль в $\R^n$, а $S_r(0) \cap \{x_n = 0\}$ имеет меру нуль как подмногообразие меньшей размерности на сфере, получаем требуемую формулу.
	Конечность внутреннего интеграла для почти всех $r$ следует из теоремы Фубини. 
\end{myproof}


\begin{corollary}
	Поверхностная мера $\sigma_{n-1}$ единичной сферы $S_1(0) = \{x \in \R^n : |x| = 1\}$ выражается через объем $\omega_n$ единичного шара $B_1(0)$ как $\sigma_{n-1} = n \omega_n$.
\end{corollary}

\begin{myproof}
	Применим теорему к функции-индикатору шара $I_{B_1(0)}(x)$:
	\[ \omega_n = \int_{\R^n} I_{B_1(0)}(x) \, d\mu(x) = \int_{0}^{+\infty} \left( \int_{S_r(0)} I_{B_1(0)}(x) \, d\nu(x) \right) \, dr. \]
	Внутренний интеграл равен площади сферы $\nu(S_r(0))$ при $r < 1$, так как в этом случае $S_r(0) \subset B_1(0)$, и нулю иначе.
	Учитывая, что $\nu(S_r(0)) = r^{n-1} \cdot \nu(S_1(0)) = r^{n-1} \sigma_{n-1}$, имеем:
	$\omega_n = \int_{0}^{1} r^{n-1} \sigma_{n-1} \, dr = \sigma_{n-1} \left[ \frac{r^n}{n} \right]_0^1 = \frac{\sigma_{n-1}}{n}$, что и требовалось. 
\end{myproof}

\begin{problem}
	Покажите, что если $x \mapsto f(|x|)$ неотрицательная измеримая или интегрируемая функция на $\R^n$, то справедлива формула 
	\[ \int_{\R^n} f(|x|) \, d\mu(x) = \sigma_{n-1} \int_{0}^{+\infty} f(r) \cdot r^{n-1} \, dr. \]
\end{problem}

\section{Дифференциальные 1-формы}

\subsection{Криволинейные интегралы второго рода}

\begin{definition}
	Пусть $U \subset \R^n$ открыто. \textit{Дифференциальной 1-формой} на $U$ называется отображение $\omega : U \to (\R^n)^*$, которое каждой точке $x \in U$ сопоставляет линейный функционал $\omega(x) : \R^n \to \R$.
\end{definition}


\begin{note}
	На 1-форму удобно смотреть как на функцию $\omega : U \times \R^n \to \R$, линейную по второму аргументу, полагая $\omega(x, h) := \omega(x)(h)$.
	Если $\{e_1, \dotsc, e_n\}$ -- стандартный базис в $\R^n$, то любой вектор $h$ разлагается как $h = \sum_{i=1}^n h_i e_i$.
	Определим \textit{коэффициенты 1-формы} $f_i : U \to \R$ как ее значения на базисных векторах: $f_i(x) = \omega(x, e_i)$.
	Тогда по линейности:
	\[ \omega(x, h) = \omega\left(x, \, \sum_{i=1}^n h_i e_i\right) = \sum_{i=1}^n h_i \, \omega(x, e_i) = \sum_{i=1}^n h_i f_i(x). \]
	Вспомним, что базис $\{dx_1, \dotsc, dx_n\}$ пространства $(\R^n)^*$, двойственный к $\{e_1, \dotsc, e_n\}$, действует как взятие $i$-й координаты: $dx_i(h) = h_i$. 
	Следовательно, имеет место \textit{координатное представление 1-формы}:
	\[ \omega(x) = f_1(x) \, dx_1 + \dotsc + f_n(x) \, dx_n. \]
	Отметим, что сложение и умножение 1-форм на скаляр определяется поточечно.
	Множество 1-форм относительно данных операций образует линейное пространство.
\end{note}

\begin{definition}
	Будем говорить, что 1-форма \textit{непрерывна}, если все ее коэффициенты непрерывны.
	Аналогично определяется принадлежность классу $C^r(U)$, $r \in \N \cup \{\infty\}$.
\end{definition}


\begin{example}
	Если $f \in C^1(U)$, то $df$ -- непрерывная 1-форма на $U$.
	Действительно, $df = \sum_{i=1}^n \frac{\partial f}{\partial x_i} dx_i$, а частные производные $\frac{\partial f}{\partial x_i}$ непрерывны по определению класса $C^1$.
\end{example}

\begin{definition}
	Пусть $\gamma : [a, b] \to U$ -- гладкая параметризованная кривая, $\omega$ -- непрерывная 1-форма на $U$. 
	Интеграл от $\omega$ по кривой $\gamma$ (\textit{криволинейный интеграл II-го рода}) определяется следующим образом: 
	\[ \int_{\gamma} \omega := \int_{a}^{b} \omega(\gamma(t), \gamma'(t)) \, dt = \int_{a}^{b} \left[ f_1(\gamma(t)) \gamma_1'(t) + \dotsc + f_n(\gamma(t)) \gamma_n'(t) \right] \, dt. \]
\end{definition}

\begin{note}
	Интеграл не зависит от параметризации. 
	Пусть $\widetilde{\gamma} : [c, d] \to U$ -- кривая, эквивалентная $\gamma$. 
	В классе гладких кривых это означает, что найдется замена параметра $h : [c, d] \to [a, b]$ -- $C^1$-сюръекция с $h' > 0$ (диффеоморфизм, сохраняющий ориентацию), такая что $\forall u \in [c, d] : \widetilde{\gamma}(u) = \gamma(h(u))$. 
	Следовательно, $\widetilde{\gamma}'(u) = \gamma'(h(u)) \cdot h'(u)$.
	Пользуясь линейностью формы по второму аргументу и формулой замены переменной ($t = h(u)$):
	\[ \int_{a}^{b} \omega(\gamma(t), \gamma'(t)) \, dt = \int_{c}^{d} \omega(\gamma(h(u)), \gamma'(h(u))) \cdot h'(u) \, du = \int_{c}^{d} \omega(\widetilde{\gamma}(u), \widetilde{\gamma}'(u)) \, du. \]
	То есть $\int_{\gamma} \omega = \int_{\widetilde{\gamma}} \omega$.
\end{note}

\begin{proposition}
	Справедливы следующие свойства интеграла: 
	\begin{enumerate}[itemsep = 0pt, topsep = 5pt]
	\item Пусть $\widetilde{\gamma} : [a, b] \to U$, $\widetilde{\gamma}(t) = \gamma(a + b - t)$, тогда $\int_{\widetilde{\gamma}} \omega = - \int_{\gamma} \omega$. %(смена ориентации).
	\item Пусть $\alpha, \beta \in \R$, тогда $\int_{\gamma} (\alpha \omega_1 + \beta \omega_2) = \alpha \int_{\gamma} \omega_1 + \beta \int_{\gamma} \omega_2$ (линейность).
	\item Пусть $a < c < b$, $\gamma_1 = \gamma|_{[a, c]}$ и $\gamma_2 = \gamma|_{[c, b]}$, тогда $\int_{\gamma} \omega = \int_{\gamma_1} \omega + \int_{\gamma_2} \omega$ (аддитивность).
	\item $\left|\int_{\gamma} \omega\right| \leq \max \limits_{x \in [\gamma]} |f(x)| \cdot L(\gamma)$, где $f = (f_1, \dotsc, f_n)$ -- вектор коэффициентов формы, $L(\gamma)$ -- длина кривой, $[\gamma]$ -- носитель кривой.
\end{enumerate}
\end{proposition}

\begin{myproof}
	Первый пункт очевиден, если вычислить интеграл с заменой $u = a + b - t$.
	Линейность следует из определения операций над формами и линейности интеграла Римана, аддитивность следует из аддитивности интеграла Римана по промежутку.
	Докажем оценку.
	Заметим, что под интегралом 1-формы по определению находится скалярное произведение $(f(\gamma(t)), \gamma'(t))$, поэтому
	\[ \left| \int_{\gamma} \omega \right| = \left| \int_a^b (f(\gamma(t)), \gamma'(t)) \, dt \right| \leq \int_a^b |(f(\gamma(t)), \gamma'(t))| \, dt \overset{\text{КБШ}}{\leq} \int_a^b |f(\gamma(t))| \cdot |\gamma'(t)| \, dt. \]
	Оценивая $|f(\gamma(t))|$ максимумом на носителе, получаем
	\[ \int_a^b |f(\gamma(t))| \cdot |\gamma'(t)| \, dt \leq \max \limits_{x \in [\gamma]} |f(x)| \int_a^b |\gamma'(t)| \, dt = \max \limits_{x \in [\gamma]} |f(x)| \cdot L(\gamma), \]
	что и требовалось.
\end{myproof}

\begin{reminder}
	Параметризованная кривая $\gamma : [a, b] \to \R^n$ называется \textit{кусочно-гладкой}, если существует такое разбиение $T = \{t_k\}_{k=0}^N$ отрезка $[a, b]$, что каждое сужение $\gamma|_{[t_{k-1}, t_k]}$ ($1 \leq k \leq N$) является гладкой параметризованной кривой. %при $k$ от 1 до $N$.
В частности, если каждое сужение $\gamma|_{[t_{k-1}, t_k]}$ является параметризованным отрезком, то $\gamma$ называется \textit{ломаной}. 
\end{reminder}

Интеграл от непрерывной 1-формы по кусочно-гладкой кривой определяется как сумма интегралов по участкам гладкости, то есть по отрезкам разбиения $T$.

\begin{theorem}
	Пусть $U \subset \R^n$ открыто, $F \in C^1(U)$ и $\gamma : [a, b] \to U$ -- кусочно-гладкая кривая. 
	Если $\gamma(a) = p$, $\gamma(b) = q$, то $\int_{\gamma} d F = F(q) - F(p)$.
\end{theorem}

\begin{myproof}
	Предположим сначала, что кривая $\gamma$ гладкая. Тогда 
	\[ \int_{\gamma} \, d F = \int_{a}^{b} \sum \limits_{i=1}^{n} \frac{\partial F}{\partial x_i} (\gamma(t)) \cdot \gamma'(t) \, dt = \int_{a}^{b} \frac{d}{dt} F(\gamma(t)) \, dt = [F(\gamma(t))]_{a}^{b} = F(q) - F(p). \]
	Для кусочно-гладкой кривой утверждение получается по аддитивности.
\end{myproof}

\begin{corollary}
	Если в условиях теоремы $\gamma : [a, b] \to U$ замкнута ($\gamma(a) = \gamma(b)$), то 
	$\int_{\gamma} \, d F = 0$.
\end{corollary}

\begin{definition}
	Пусть $U \subset \R^n$ открыто, $\omega$ -- 1-форма в $U$. 
	\begin{enumerate}[itemsep = 0pt, topsep = 5pt]
		\item $F : U \to \R$ называется \textit{первообразной} $\omega$, если $F$ дифференцируема и $d F = \omega$ на $U$, 
		\item Форма $\omega$ называется \textit{точной в $U$}, если она имеет там первообразную.
	\end{enumerate}
\end{definition}

В дальнейшем нам потребуется один факт относительно связных открытых множеств.

\begin{lemma}
	Если $U$ -- область в $\R^n$, то любую пару точек из $U$ можно соединить ломаной со сторонами, параллельными осям координат.
\end{lemma}


\begin{myproof}
	Внутри открытого шара $B = B_r(a)$ центр $a$ можно соединить с любой точкой $y \in B$ искомой ломаной.
	Действительно, пусть $a = (a_1, \dots, a_n)$ и $y = (y_1, \dots, y_n)$, построим последовательность $z_0, \dots, z_n$, меняя координаты $a$ на координаты $y$ по одной:
	\[ z_0 = (a_1, a_2, \dotsc, a_n), \, z_1 = (y_1, a_2, \dots, a_n), \, \dotsc, \, z_k = (y_1, \dots, y_k, a_{k+1}, \dots, a_n). \]
	Отрезок $[z_{k-1}, z_k]$ параллелен $k$-й оси координат. Проверим, что все эти точки и отрезки лежат внутри шара.
	Расстояние от центра $a$ до промежуточной точки $z_k$:
	\[ \rho(a, z_k) = \sqrt{\sum_{i=1}^k (y_i - a_i)^2 + \sum_{j=k+1}^n (a_j - a_j)^2} = \sqrt{\sum_{i=1}^k (y_i - a_i)^2} \le \sqrt{\sum_{i=1}^n (y_i - a_i)^2} = \rho(a, y) < r. \]
	Так как $z_{k-1}, z_k \in B_r(a)$, а шар выпуклый, то и весь отрезок $[z_{k-1}, z_k] \subset B_r(a)$.
	
	\vspace{5pt}
	Докажем для $U$. Зафиксируем $x_0 \in U$. Рассмотрим множество $A$ всех точек $x \in U$, достижимых из $x_0$ ломаной со сторонами, параллельными осям.
	Множество $A \neq \emptyset$, так как $x_0 \in A$.
	Пусть $x \in A$. Так как $U$ открыто, то $\exists B_r(x) \subset U$. Любую точку $y \in B_r(x)$ можно соединить с центром $x$ ломаной внутри шара. Объединяя пути $x_0 \rightsquigarrow x$ и $x \rightsquigarrow y$, получаем $y \in A$. Следовательно, $B_r(x) \subset A$ и $A$ открыто.
	Пусть $y \in U \setminus A$. Аналогично $\exists B_r(y) \subset U$. Если бы в этом шаре нашлась достижимая точка $x \in A$, то мы могли бы пройти по пути $x_0 \rightsquigarrow x \rightsquigarrow y$, соединив $x$ и $y$ внутри шара, что противоречит недостижимости $y$. Следовательно, $B_r(y) \subset U \setminus A$ и $U \setminus A$ открыто.
	В силу связности $U$, единственным кандидатом на множество $A$ является $U$. Значит, $A=U$, что и требовалось.
\end{myproof}

\begin{theorem}
	Пусть $U$ -- область в $\R^n$ и $\omega$ -- непрерывная 1-форма на $U$. Следующие утверждения эквивалентны: 
	\begin{enumerate}[itemsep = 0pt, topsep = 5pt]
		\item Форма $\omega$ точна в $U$, 
		\item $\int_{\gamma} \omega = 0$ по любой замкнутой кусочно-гладкой кривой $\gamma$ из $U$, 
		\item $\int_{\gamma_1} \omega = \int_{\gamma_2} \omega$ для любых ломаных $\gamma_i : [a, b] \to U$ с совпадающими концами и со сторонами, параллельными осям координат.
	\end{enumerate}
\end{theorem}

\begin{myproof}
	$(1 \Ra 2)$ Вытекает из следствия первой теоремы. 

	\vspace{5pt}
	$(2 \Ra 3)$ Зафиксируем ломаные $\gamma_i$ с общими концами и рассмотрим кривую $\gamma$: 
	\[ \gamma(t) := \begin{cases}
	 \gamma_1(t), & t\in [a, b],   \\  
	 \gamma_2(2b - t),  & t\in [b, 2b-a].  
	\end{cases} \] Так как $\gamma$ -- замкнутая кусочно-гладкая кривая, то $0 = \int_{\gamma} \omega = \int_{\gamma_1} \omega - \int_{\gamma_2} \omega$.

	\vspace{5pt}
	$(3 \Rightarrow 1)$ Зафиксируем $x_0 \in U$. Рассмотрим функцию $F(x) = \int_{\gamma_{x_0, x}} \omega$, где интеграл берется по ломаной $\gamma_{x_0, x}$ со сторонами, параллельными осям, соединяющей $x_0$ и $x$. Существование ломаной доказано в лемме. 
	Согласно пункту (3), значение $F(x)$ не зависит от выбора конкретной ломаной. 
	Покажем, что $F$ -- это первообразная для $\omega$. Пусть $\omega = \sum_{k=1}^{n} f_k(x) dx_k$, тогда достаточно проверить, что $\forall i : \frac{\partial F}{\partial x_i} = f_i$ на $U$.
	Пусть $x \in U$, тогда $\exists \delta > 0$, что $\forall t \in (-\delta, \delta)$ точка $x + t e_i \in U$.
	Параметризуем отрезок, соединяющий $x$ и $x + t e_i$: $\lambda_i(s) = x + s e_i$.
	Найдем приращение функции, пользуясь независимостью от пути:
	\[ F(x + t e_i) - F(x) = \int_{\gamma_{x_0, x+te_i}} \omega - \int_{\gamma_{x_0, x}} \omega = \int_{\lambda_i} \omega = \int_0^t \omega(x+se_i, e_i) \, ds = \int_0^t f_i(x + se_i) \, ds. \]
	Тогда мы можем найти частную производную по определению:
	\[ \frac{\partial F}{\partial x_i}(x) = \lim \limits_{t \to 0} \frac{F(x + t e_i) - F(x)}{t} = \lim \limits_{t \to 0} \frac{1}{t} \int_0^t f_i(x + se_i) \, ds = \frac{d}{dt} \Big|_{t = 0} \int_0^t f_i(x + se_i) \, ds \]
	Производная интеграла с переменным верхним пределом равна $f_i(x)$.
 \end{myproof}

 \begin{corollary}[\textbf{лемма Гурса}]
	При $n=2$ для точности формы достаточно проверять равенство нулю интеграла по любому прямоугольнику со сторонами, параллельными осям. 
 \end{corollary}

\begin{myproof}
	Любую замкнутую ломаную со сторонами, параллельными осям, можно представить как границу объединения конечного числа прямоугольников. 
	В силу аддитивности, если интеграл по границе каждого прямоугольника равен нулю, то интеграл по исходному контуру также равен нулю. 
	Это сводит утверждение к третьему пункту теоремы.
\end{myproof}
